\section{Polynomial expected depth with adaptive block modes}
\label{sec:adaptive-blocks}

Sometimes querying the second endpoint of a block remains useful after
the first endpoint was positive. The additional answer can pin an arc
boundary and reveals information that would otherwise be learned by
endpoint decoding. Accounting for both effects, while choosing every
mode by rational arithmetic, improves the general expected-depth factor
to below $5.272565$. The proof has two independent components:
a local two-endpoint saving, proved below, and a finite coverage-prefix
certificate. Lemmas~\ref{lem:capped-rank-credit}
and~\ref{lem:stage-conversion} then convert that certificate into
the full recursive guarantee.

\subsection{Physical costs and a coarsened transcript}

Keep the polynomial exact-coverage and geometric-batch construction of
Section~\ref{sec:polynomial-blocks}. In a block of rank $k$, let $x,y$
be its newly hit masses exclusive to its two endpoints and $z$ their
shared mass. Orient the block so $x\le y$, and put $m=x+y+z$.
Each endpoint cap has mass at most $\delta w$, where $w=w(S)$;
thus $x+z,y+z\le\delta w$.

In \emph{stop mode}, stop querying this block at the first positive
answer. Relative to the nonentropy allowance $(2k+3)m$, the $y$ mass
exits second with a pinned boundary and saves one query per hypothesis.
The saving is $y$.

In \emph{query-both mode}, ask the second endpoint even when the first
was positive. Each exclusive mass has one adjacent negative endpoint,
so only one alphabetic endpoint code is needed. Its nonentropy cost
is at most $2k+2$. The shared mass has two positive endpoints and
costs at most $2k+4$. The physical saving is therefore $x+y-z$.
An origin endpoint is free, and cached answers only lower these costs.
For the analysis, complete the chosen mode before singleton inference;
constant queries can be contracted afterwards.

Let $\mathcal R_{\rm act}$ be the actual discovery partition and
$J_{\rm act}=I(S,\mathcal R_{\rm act})$. Define a coarsened
partition $\mathcal R_*$ by forgetting the second answer on each
query-both route whose first answer was positive. This merges that
block's $x$ and $z$ events. Its $y$ event remains separate. Stop-mode
events are unchanged. Every nonzero coarsened event lies in a single
positive endpoint cap, so, for $J_*=I(S,\mathcal R_*)$,
\begin{equation}\label{eq:coarse-discovery-cap}
 J_*\ge b\sigma(\delta),\qquad b=w-p_0.
\end{equation}
The zero event is the same zero-pattern singleton as before.

The coarsening loses exactly
\[
 J_{\rm act}-J_*
 =\sum_{\text{query-both blocks}}
 (x+z)h\!\left(\frac{x}{x+z}\right),
\]
with zero contribution when $x+z=0$. Here
$h(p)=-p\ln p-(1-p)\ln(1-p)$. Actual endpoint decoding costs
at most its physical allowance plus $(I-J_{\rm act})/\ln2$,
where $I=I(S,\Pi)$ is the full-pattern information. Rewriting this
term as $(I-J_*)/\ln2$ therefore credits an additional
$(x+z)H_2(x/(x+z))$ queries on every query-both block. The useful ledger is therefore physical saving plus the information
forgotten by coarsening. The remaining-information term uses $J_*$
throughout, so that information is charged exactly once.

\subsection{Rational mode choices and affine saving bounds}

Fix
\[
 p_* =\frac{21317}{125000},\qquad
 h_* =1-2p_* =\frac{41183}{62500}.
\]
Exact logarithm enclosures certify $H_2(p_*)>h_*$. Concavity and
symmetry of binary entropy give the rational tent bound
\[
 H_2(p)\ge\min\left\{\frac{h_*}{p_*}p,\ h_*,
                         \frac{h_*}{p_*}(1-p)\right\}.
\]
Its homogeneous mass version and the rational block saving are
\begin{equation}\label{eq:rational-block-saving}
 \begin{aligned}
 E_*(x,z)&=\min\left\{\frac{h_*}{p_*}x,\ h_*(x+z),
                                  \frac{h_*}{p_*}z\right\},\\
 D(x,y,z)&=\max\{y,\ x+y-z+E_*(x,z)\}.
 \end{aligned}
\end{equation}
Algorithm~\ref{alg:block-mode} selects the mode attaining this maximum;
Algorithm~\ref{alg:expected-tree} gives the full recursion. Both its choice and its saving
are rational functions of the supplied prior weights. Let
$\operatorname{Save}(G)=\sum_{\text{blocks of }G}D(x,y,z)$.
For each candidate geometric shift, choose its modes and minimize
the rational surrogate
\begin{equation}\label{eq:adaptive-shift-objective}
 F_*(G)=2\MSSC_{\rm blk}(G)+2p_0|G|-\operatorname{Save}(G).
\end{equation}
The full stage then obeys
\begin{equation}\label{eq:adaptive-stage}
 W_{\rm pat}(S)\le F_*(G)+3b+\frac{I-J_*}{\ln2}.
\end{equation}
Modes do not change a block's union, newly covered mass, coverage sets,
or greedy order within a batch. The shift-average coverage proof
therefore remains available, and every shift satisfies
\eqref{eq:coarse-discovery-cap}.

\begin{algorithm}[H]
\caption{Rational orientation and adaptive block mode}
\label{alg:block-mode}
\KwIn{Newly hit exclusive masses $x,y$ and shared mass $z$; $p_*,h_*$ from Section~\ref{sec:adaptive-blocks}.}
\KwOut{Endpoint order, mode, and certified rational saving $D$.}
Orient the block so $x\le y$\;
$E_*\gets\min\{(h_*/p_*)x,\ h_*(x+z),\ (h_*/p_*)z\}$\;
\eIf{$y\ge x+y-z+E_*$}{mode $\gets$ stop; $D\gets y$\;}{
mode $\gets$ query-both; $D\gets x+y-z+E_*$\;
}
Optionally compare query-both in the reversed orientation using $E_*(y,z)$,
and retain it if its saving is larger\;
Record this mode for the first hitting block; in query-both mode its
stored stage learns both endpoint bits before endpoint decoding\;
\KwRet the chosen order, mode, and $D$\;
\end{algorithm}

The new masses are calculated on the all-negative prefix before this
block. Mode choice leaves its union and first-cover rank unchanged.
The coarsened-transcript proof in Section~\ref{sec:adaptive-blocks}
accounts for the credited endpoint information exactly once.


\begin{lemma}[Affine block credits]\label{lem:affine-block-credit}
Put
\[
 c_* =\frac{228683}{103683}=\frac{2-p_*}{1-p_*},\qquad
 d_* =\frac{250000}{103683}=\frac2{1-p_*}.
\]
Every block with newly hit mass $m$ satisfies
\[
 D\ge0,\qquad D\ge m-\delta w,\qquad
 D\ge c_*m-d_*\delta w.
\]
\end{lemma}
\begin{proof}
The first two follow from $D\ge y\ge0$ and
$m-y=x+z\le\delta w$. For the third, normalize $\delta w=1$.
If $m\le1$, its right side is negative. Otherwise $1\le m\le2$,
and the feasible triangle for $x\le y$, with $z=m-x-y$, has vertices
\[
 (m-1,m-1),\qquad(m-1,1),\qquad(m/2,m/2).
\]
The query-both expression $x+y-z+E_*(x,z)$ is concave on this
triangle, hence bounded below by its smallest vertex value. The
two vertices with $z=0$ have value $m$. At the first vertex it is
\[
 f(m)=3m-4+
       \min\left\{\frac{h_*}{p_*}(m-1),\ h_*,
                            \frac{h_*}{p_*}(2-m)\right\}.
\]
For $m\le1+p_*$, $c_*m-d_*\le m-1$, already covered by the
second credit. For $m\ge1+p_*$, $f$ is concave with
$f(1+p_*)=p_*$ and $f(2)=2$. Its chord on this interval is
exactly $c_*m-d_*$. Also $m\ge c_*m-d_*$ for $m\le2$.
Both possible lower vertex values therefore dominate the required
line. Rescale by $\delta w$.
\end{proof}

\subsection{A certified mixture of coverage-prefix credits}

Keep the witness masses $m_t$ and $M_k=\sum_{t\le k}m_t$ from
Section~\ref{sec:polynomial-blocks}, with $m_t\le2\delta w$.
A comparison edge sequence covers at least $M_k$ in its first $k$
blocks. Exact budgeted coverage therefore gives covered mass at least
$M_k$ in a batch with budget $k$. If the schedule prefix through
that batch uses at most $\ell$ blocks, Lemma~\ref{lem:affine-block-credit}
and nonnegative later savings imply each of the lower bounds
\[
 \operatorname{Save}(G)\ge
 \mathsf A(k,\ell):=M_k-\delta w\ell,
 \qquad
 \operatorname{Save}(G)\ge
 \mathsf C(k,\ell):=c_*M_k-d_*\delta w\ell.
\]
It is also at least zero. Early completion only improves these bounds.
For a shift $u\in[1,\beta)$, the first batch has
$k=\lfloor u\rfloor$, $\ell\le\lfloor u\rfloor$; through the
second batch, $k=\lfloor\beta u\rfloor$ and
$\ell\le\lfloor u\rfloor+\lfloor\beta u\rfloor$.
Any convex combination of their lower bounds and zero bounds the
\emph{same} total saving, counting it once.

Use
\[
 \beta=\frac{3512863}{10^6},\qquad
 \delta=\frac{144136826}{10^9},\qquad
 v_1=\frac{830027}{10^7},\quad
 v_2=\frac{472577}{10^6},\quad
 v_3=\frac{1701301}{5\cdot10^6}.
\]
The following fixed choices are used only in the proof:
\begin{center}
\small
\begin{tabular}{ll}
\hline
Shift range & Lower bound on $\operatorname{Save}(G_u)$\\
\hline
$[1,4/\beta)$ & $\mathsf C(1,1)$\\
$[4/\beta,5/\beta)$ & $v_1\mathsf C(1,1)+(1-v_1)\mathsf A(4,5)$\\
$[5/\beta,6/\beta)$ & $\mathsf A(5,6)$\\
$[6/\beta,2)$ & $\mathsf C(1,1)$\\
$[2,10/\beta)$ & $\mathsf A(2,2)$\\
$[10/\beta,3)$ & $v_2\mathsf A(2,2)+(1-v_2)0$\\
$[3,11/\beta)$ & $v_3\mathsf A(3,3)+(1-v_3)\mathsf C(3,3)$\\
$[11/\beta,12/\beta)$ & $\mathsf C(3,3)$\\
$[12/\beta,\beta)$ & $\mathsf A(3,3)$\\
\hline
\end{tabular}
\end{center}

Here is an explicit coefficient definition for the certificate. Start
with $a_t=e_\beta(t)$ and $\mu=0$. A table row with range $[l,r)$
has log-uniform probability $P=\ln(r/l)/\ln\beta$. For each
minorant $\alpha M_k-\rho\delta w$ in that row, with mixing weight
$v'$, add $Pv'\alpha$ to every $a_t$ with $t\le k$ and add
$Pv'\rho$ to $\mu$. For a zero term both additions are zero.
No table entry has $k>5$, so $a_t=e_\beta(t)$ for $t\ge6$.
This defines fixed analytical constants without an optimization oracle.
The shift-average transfer is
\[
 \sum_t m_t e_\beta(t)+\mathbb E[\operatorname{Save}(G_u)]
 \ge\sum_t m_t a_t-\delta w\mu.
\]

Set $\tau=345603/250000=1.382412$. Exact rational logarithm
enclosures certify that only $a_4$ lies below $\tau$:
\[
 \begin{aligned}
 a_1&>1.38241213,& a_2&>1.38241211,\\
 a_3&>1.38241211,& a_4&<1.37612333,\\
 a_5&>1.38241212,& e_\beta(6)&>1.38382029.
 \end{aligned}
\]
Monotonicity of $e_\beta$ controls the remaining tail. Define
\begin{equation}\label{eq:adaptive-chi-B}
 \chi_* =\tau-2\delta(\tau-a_4)-\delta\mu,\qquad
 B_* =3+2\Phi_\beta(1)-\chi_*
      =7-\frac{\ln6}{\ln\beta}-\chi_*.
\end{equation}
This is Lemma~\ref{lem:capped-rank-credit} with the single
exceptional index $E=\{4\}$ and $\zeta=2/\ln\beta$.
The certified inequality $\chi_*+2/\ln\beta>\tau$ is its
zero-route condition. Thus the combined discrete, block-mode, and
zero-route discounts save at least $\chi_*b$.
Minimizing \eqref{eq:adaptive-shift-objective} allows the saving to
be retained under the same shift average as the first-hit and length
terms. The proof of \eqref{eq:poly-block-discrete} now gives
\begin{equation}\label{eq:adaptive-local}
 W_{\rm pat}(S)\le A\Delta+B_*b+\frac{I-J_*}{\ln2},
 \qquad A=\frac{\beta+1}{\ln\beta}.
\end{equation}
The at-most-two-real-tests degenerate-hull construction also satisfies
this bound because $A,B_*>2$, exactly as in the earlier block proof.

\begin{theorem}[Adaptive-mode polynomial expected-depth approximation]
\label{thm:poly-adaptive-expected}
For general supplied binary unit-cost planar line queries and arbitrary
strictly positive rational priors, the adaptive-mode block algorithm
runs in deterministic polynomial time and has approximation factor
\[
 \boxed{\rho_*=
 A+\max\{\ln2/h(\delta),\ B_*\ln2/\sigma(\delta)\}<5.272565.}
\]
Its exponent has no dependence on the smallest prior.
\end{theorem}
\begin{proof}
Apply Lemma~\ref{lem:stage-conversion} with $B=B_*$ to
\eqref{eq:adaptive-local} and \eqref{eq:coarse-discovery-cap}.
The coarsened transcript is a function of the full pattern, so
$0\le J_*\le I$ and the lemma applies unchanged.

The exact rational certificate \texttt{further-expected-check.py}
verifies the tent inequality, all coefficient comparisons, the
zero-route condition, and
\[
 \begin{aligned}
 0.856049452254&<\chi_*<0.856049452257,\\
 4.717880246117&<B_*<4.717880246121,\\
 \max\{\ln2/h(\delta),B_*\ln2/\sigma(\delta)\}
 &<1.680753968384.
 \end{aligned}
\]
Together with $A<3.591810178840040$, these certify
$\rho_*<5.272564147225<5.272565$.
The script also checks 5,000 rational block-credit instances and
10,000 analytical source profiles; its theorem certificate does not
use floating-point optimization.

Coverage DPs, rational shift enumeration, and the disjoint-pattern
recursion remain polynomial. Each mode uses fixed rational linear
forms in the prior weights. Computing the entropy, logarithms, or
analysis-table mixture is unnecessary for constructing the tree.
Every query remains a supplied endpoint line.
\end{proof}

\paragraph{Further refinements.}
For an instance-level improvement, query-both mode can also use the
larger exclusive mass first, replacing $E_*(x,z)$ by $E_*(y,z)$.
This is no smaller and its coarsened exits retain the same cap.
Choosing the best of both orientations and stop mode by the rational
surrogate preserves the theorem. A larger constant improvement would
require additional scheduling or accounting structure; the present
affine proof is tight at symmetric exclusive masses and cannot alone
benefit from this orientation change.
